\documentclass[aspectratio=169,10pt,spanish]{beamer}

\input{../../../_preambulo_beamer.tex}
\input{../../../_comandos_beamer.tex}

\selectlanguage{spanish}

\title{MA1001B - Modelación Estadística para la Toma de Decisiones}
\subtitle{Lección 42: Prueba ji-cuadrada de homogeneidad}
\author{}
\institute{Tecnol\'ogico de Monterrey}
\date{}

\begin{document}

\begin{frame}[plain]
  \vspace{1.2cm}
  {\Large\bfseries MA1001B - Modelación Estadística para la Toma de Decisiones\par}
  \vspace{0.7cm}
  {\large Lección 42: Prueba ji-cuadrada de homogeneidad\par}
  \vspace{0.7cm}
  {\color{TecRojo}\rule{\textwidth}{1.2pt}}
  \vspace{0.8cm}

  Facultad MA1001B\\[0.3cm]
  Periodo\\[0.3cm]
  Tecnol\'ogico de Monterrey
\end{frame}

\begin{frame}{La pregunta de la mezcla compartida}
  \begin{alertblock}{Pregunta guía}
    ¿Comparten tres almacenes muestreados de forma independiente una mezcla
    de canal de tickets?
  \end{alertblock}

  \vspace{0.6em}
  Al final de esta lección, usted puede:
  \begin{itemize}
    \item tratar los totales de fila como fijos por el plan de muestreo;
    \item calcular una $\chi^2$ de homogeneidad con $df=(r-1)(c-1)$;
    \item decidir si los sitios comparten una mezcla;
    \item contrastar esa historia con una prueba de independencia.
  \end{itemize}

  \vfill
  \scriptsize Todos los conteos de tickets usados hoy son sintéticos. No se usan datos reales de empresa.
\end{frame}

\begin{frame}{Tres muestras planeadas, mismas categorías}
  \small
  \begin{center}
    \begin{tabular}{lrrrr}
      \toprule
      \textbf{Sitio} & \textbf{Teléfono} & \textbf{Chat} & \textbf{Correo} & \textbf{$n$} \\
      \midrule
      Norte & 60 & 40 & 20 & 120 \\
      Centro & 35 & 40 & 25 & 100 \\
      Sur & 25 & 30 & 35 & 90 \\
      \midrule
      Total & 120 & 110 & 80 & 310 \\
      \bottomrule
    \end{tabular}
  \end{center}

  \vspace{0.3em}
  \begin{exampleblock}{Hipótesis de homogeneidad}
    $H_0$: los tres sitios comparten una mezcla frente a $H_1$: al menos un
    sitio difiere, con $\alpha=0.05$.
  \end{exampleblock}
\end{frame}

\begin{frame}[fragile]{Paso 1: Tres mezclas independientes}
  \begin{columns}[T]
    \begin{column}{0.42\textwidth}
      \small
      \begin{lstlisting}
shares = counts / n_site
      \end{lstlisting}

      \begin{block}{Salida verificada}
        Participación telefónica Norte $=0.500000$\\
        Participación telefónica Centro $=0.350000$\\
        Participación telefónica Sur $=0.277778$
      \end{block}
    \end{column}
    \begin{column}{0.58\textwidth}
      \centering
      \includegraphics[width=\textwidth]{../figuras/homogeneidad_01_tres_muestras.png}
      \vspace{0.2em}
      \scriptsize Mezclas dentro del sitio del Paso 1
    \end{column}
  \end{columns}
\end{frame}

\begin{frame}{La homogeneidad usa la misma aritmética $\chi^2$}
  \small
  Los conteos esperados siguen usando el producto de márgenes:
  \[
    E_{ij}=\frac{(\text{total de fila}_i)(\text{total de columna}_j)}{n}.
  \]
  \[
    \chi^2=17.812605,\qquad df=(3-1)(3-1)=4,\qquad p=0.001343.
  \]
  Como $\chi^2>\chi^{2*}=9.487729$ y $p<\alpha$, la prueba rechaza la
  afirmación de que los tres sitios comparten una mezcla.
\end{frame}

\begin{frame}[fragile]{Paso 2: Decisión de homogeneidad}
  \begin{columns}[T]
    \begin{column}{0.40\textwidth}
      \small
      \begin{lstlisting}
chi2, p, df, E = chi2_contingency(table)
      \end{lstlisting}

      \begin{block}{Salida verificada}
        $\chi^2=17.812605$, $df=4$\\
        $p=0.001343$\\
        Decisión: rechazar $H_0$
      \end{block}
    \end{column}
    \begin{column}{0.60\textwidth}
      \centering
      \includegraphics[width=\textwidth]{../figuras/homogeneidad_02_chi_cuadrada.png}
      \vspace{0.2em}
      \scriptsize Contribuciones celulares de homogeneidad del Paso 2
    \end{column}
  \end{columns}
\end{frame}

\begin{frame}{Paso 3: El mismo estadístico, distinta historia}
  \begin{columns}[T]
    \begin{column}{0.42\textwidth}
      \small
      Homogeneidad: tres muestras planeadas. Los totales de fila $120$,
      $100$ y $90$ se eligieron primero.

      \vspace{0.4em}
      Independencia: una muestra de $310$ tickets clasificada después por
      sitio y canal.

      \vspace{0.5em}
      \textbf{Límite:} ambas historias producen $\chi^2=17.812605$ y
      $p=0.001343$. La aritmética se comparte; la $H_0$ no.
    \end{column}
    \begin{column}{0.58\textwidth}
      \centering
      \includegraphics[width=\textwidth]{../figuras/homogeneidad_03_limite_historia_muestreo.png}
      \vspace{0.2em}
      \scriptsize Dos historias de muestreo del Paso 3
    \end{column}
  \end{columns}
\end{frame}

\begin{frame}{Verificación de conceptos}
  \begin{enumerate}
    \item ¿Por qué los totales de fila $120$, $100$ y $90$ son parte del
      diseño?
    \item Calcule $df$ para tres sitios y tres categorías.
    \item Si $p=0.001343$ y $\alpha=0.05$, ¿cuál es la decisión de
      homogeneidad?
    \item ¿Cómo puede la misma tabla ser un problema de independencia?
  \end{enumerate}

  \vspace{0.6em}
  \begin{block}{Conexión con la Lección 43}
    La sesión puede continuar directamente con diseño experimental:
    factores, aleatorización y por qué un sitio observacional no es un
    tratamiento.
  \end{block}

  \vfill
  \scriptsize Fuente: OpenStax, \textit{Introductory Business Statistics 2e},
  Capítulo 11, Sección 11.5. CC BY-NC-SA.
\end{frame}

\appendix



\end{document}
